% Part: incompleteness
% Chapter: introduction
% Section: decidability

\documentclass[../../../include/open-logic-section]{subfiles}

\begin{document}

\olfileid{inc}{int}{dec}

\olsection{不可判定性与不完全性}

哥德尔 对不完全性定理的证明需要对语法进行算术化。
但即使没有这一点，仅基于一个理论能表示所有可判定关系这一假设，我们也能得到一些很好的结果。
其证明是一个类似于停机问题不可判定性证明的对角线论证。

\begin{thm}
若 $\Gamma$ 是一个表示所有可判定关系的一致理论，则 $\Gamma$ 不可判定。
\end{thm}

\begin{proof}
假设 $\Gamma$ 是可判定的。我们证明，若 $\Gamma$
表示所有可判定关系，则它必然不一致。

可判定的性质（一元关系）由含一个自由变元的公式表示。
令 $!A_0(x)$, $!A_1(x)$, \dots,
为所有这些公式的一个可计算枚举。现在考虑以下集合 $D \subseteq \Nat$：
\[
D = \Setabs{n}{\Gamma \Proves \lnot !A_n(\num{n})}
\]
集合 $D$ 是可判定的，因为我们可以先计算 $!A_n(x)$，再由其 $\lnot !A_n(\num{n})$，从而检验是否有 $n \in D$。显然，
将项 $\num{n}$ 代入 $!A_n(x)$ 中 $x$ 的所有自由出现，
并在 $!A(\num{n})$ 前加上 $\lnot$，这些都是机械操作。
根据假设，$\Gamma$ 是可判定的，因此我们可以检验是否
$\lnot !A(\num{n}) \in \Gamma$。若是，则 $n \in D$；若不是，则
$n \notin D$。所以 $D$ 同样是可判定的。

由于 $\Gamma$ 表示所有可判定性质，它
也表示~$D$。而在~$\Gamma$ 中表示 $D$
的公式都在 $!A_0(x)$, $!A_1(x)$, \dots 之中。
令 $d$ 是满足 $!A_d(x)$ 在~$\Gamma$ 中表示 $D$ 的数。
若 $d \notin D$，那么由于 $!A_d(x)$ 表示~$D$，$\Gamma \Proves
\lnot !A_d(\num{d})$。但这意味着 $d$ 满足~$D$ 的定义，
因此 $d \in D$。这与 $d \notin D$ 矛盾。
故由反证法，$d \in D$。

既然 $d \in D$，由~$D$ 的定义，$\Gamma \Proves \lnot
!A_d(\num{d})$。另一方面，由于 $!A_d(x)$ 在 $\Gamma$ 中表示~$D$，
$\Gamma \Proves !A_d(\num{d})$。因此，$\Gamma$ 是不一致的。
\end{proof}

\begin{explain}
前述定理表明，任何表示所有可判定关系的一致理论都不可能是可判定的。
我们将证明 $\Th{Q}$ 确实表示了所有可判定
关系；这意味着所有包含 $\Th{Q}$ 的理论，例如
$\Th{PA}$ 和 $\Th{TA}$，也都表示所有可判定关系，因而也都不可判定。
（由于这些理论在标准模型中为真，
它们都是一致的。）

我们也可以利用这个结果来得到第一不完全性定理的一个弱形式。
任何可公理化且完全的理论都是可判定的。
因此，一致、可公理化且表示所有可判定性质的理论
就不可能是完全的。
\end{explain}

\begin{thm}
若 $\Gamma$ 可公理化且完全，则它是可判定的。
\end{thm}

\begin{proof}
任何不一致的理论都是可判定的，因为不一致的理论
包含所有语句，所以对“$!A \in
\Gamma$ 吗？”这一问题的回答总是“是”，即可判定。

所以假设 $\Gamma$ 是一致的，并且是
可公理化和完全的。由于 $\Gamma$ 可公理化，
它是可计算枚举的。因为我们可以用一个可计算函数
枚举从 $\Gamma$ 的公理出发的所有正确推导。
由一个正确推导可以计算出它所推导的语句，
因此综合起来，存在一个可计算函数能枚举 $\Gamma$ 的所有定理。
由于 $\Gamma$ 是一致且完全的，语句 $!A$ 是 $\Gamma$ 的定理，当且仅当 $\lnot !A$ 不是定理。
因此，我们可以如下判定是否 $!A \in
\Gamma$：枚举 $\Gamma$ 的所有定理。当 $!A$
出现在此列表中时，可知 $\Gamma \Proves !A$。当 $\lnot
!A$ 出现在此列表中时，可知 $\Gamma \Proves/ !A$。由于
$\Gamma$ 是完全的，这两种情况最终必居其一，因此
该程序最终必能给出答案。
\end{proof}

\begin{cor}
\ollabel{cor:incompleteness}
若 $\Gamma$ 是一致的、可公理化的，并且表示所有
可判定性质，则它不是完全的。
\end{cor}

\begin{proof}
若 $\Gamma$ 是完全的，则根据前一定理（因为 $\Gamma$ 可公理化且一致），
它是可判定的。但既然 $\Gamma$ 表示所有可判定性质，由前一定理，
它是不可判定的。
\end{proof}

\begin{prob}
证明 $\Th{TA} = \Setabs{!A}{\Sat{N}{!A}}$ 不可公理化。可假设 $\Th{TA}$ 表示所有可判定性质。
\end{prob}

一旦确立了例如 $\Th{Q}$ 表示所有可判定性质，
该推论便告诉我们 $\Th{Q}$ 必定是不完全的。
然而，该推论的证明并未给出独立语句的实例；
它仅仅表明这样的语句必然存在。
为此，我们必须对语法进行算术化，并遵循 哥德尔 的原始证明思路。
当然，我们还得证明第一个断言，即 $\Th{Q}$ 确实表示了所有可判定性质。

应当注意的是，并非每个\emph{有趣的}理论都是不完全的或不可判定的。
有许多理论足以描述有趣的数学事实，却并不满足 哥德尔 结果的条件。
例如，$\Th{Pres} = \Setabs{!A \in \Lang{L_{A^+}}}{\Sat{N}{!A}}$（即在标准模型中为真、不含~$\times$ 的算术语言语句的集合）既是完全的，又是可判定的。
这个理论称为 Presburger 算术，它证明了所有仅用 $\Obj{0}$、$\prime$ 和~$+$ 就能表述的关于自然数的真命题。

\end{document}